\section{Intuitive Picture}


直観的には、システム内には少なくとも二つの速度がある。ひとつは、予測、反応、瞬間的更新を担う速い流れである。もうひとつは、記憶、安定化、長期的構造変化を担う遅い流れである。

意識は速い流れそのものではない。また、遅い記憶層そのものでもない。意識は、速い流れが遅い層に完全に沈殿する直前、あるいはその沈殿が可能になる境界で開く。

この境界は静止した場所ではない。高速過程と低速過程の速度差が維持され、かつ結合が切れないときにだけ成立する動的な関係である。速度差が消えれば境界はなくなる。速度差が大きすぎれば参照は破断する。
